% Exact audit of the blow-up route; owned by cycles. No preamble.
\paragraph{C.1q (Proved: the exact exceptional polynomial and the ample-pencil test).}
Let $M/\mathbb C$ be smooth projective of dimension $2m$, let
$i:Z\hookrightarrow M$ be a smooth connected center of codimension
$c\ge2$, and put $\pi:X=\operatorname{Bl}_Z(M)\to M$.
Write $j:E\hookrightarrow X$ and $p:E\to Z$ for the exceptional
divisor and its projection, and set
\[
 N=N_{Z/M},\qquad \xi=c_1(\mathcal O_{\mathbb P_{\mathrm{lines}}(N)}(1)).
\]
Assume $G^m(M)=G^{m+1}(M)=0$; in particular this holds when
$M$ has a Tate Chow motive. Let
$L=\pi^*D-bE$ be an ample rational class, so $b>0$, and put
$u=i^*D/b$. Define
\[
 q_N(T)=\sum_{k=0}^{c-1}c_k(N)T^{c-1-k},\qquad c_0(N)=1.
\]
Then extraction of the last exceptional coefficient gives an
isomorphism
\begin{equation}\label{eq:C1q-kernel}
 \operatorname{rad}(q_{X,L})\simeq
 \ker\left(q_N(-u):G^{m-c+1}(Z)\longrightarrow G^m(Z)\right).
\end{equation}
Thus Griffiths classes strictly above the middle degree of a
center, by themselves, do not create a primitive middle radical:
the controlling degree is
\[
             m-c+1=\frac{\dim Z-c+2}{2}\le\frac{\dim Z}{2}.
\]
No condition on the Chern classes of the normal bundle can change
that degree constraint.

\begin{proof}
We use the opened blow-up formula
\citep[Theorem 2.8(3), and equation (2.2)]{cyJiang2023}.
Its maps and inverse are algebraic correspondences, so the
formula descends to homological equivalence and to algebraic
equivalence, and hence to the rational Griffiths quotient.
Every $\Gamma\in G^m(X)$ has a unique expression
\[
 \Gamma=\sum_{j=1}^{c-1}f_j(x_j),\qquad
 f_j(x)=j_*(\xi^{j-1}p^*x),\quad x_j\in G^{m-j}(Z).
\]
The excess bundle is $Q=p^*N/\mathcal O_E(-1)$. Its top Chern
class is
\[
 c_{c-1}(Q)=\sum_{k=0}^{c-1}p^*c_k(N)\xi^{c-1-k}.
\]
The excess identity therefore gives
\[
 f_c(z)=\pi^*i_*z-\sum_{k=1}^{c-1}f_{c-k}(c_k(N)z).
\]
For the class $z$ used below, the ambient term lies in
$G^{m+1}(M)$ and is zero. Also $j^*E=-\xi$, so
\[
 b^{-1}L\,f_j(x)=f_j(ux)+f_{j+1}(x).
\]
These formulas fix the signs and include the top exceptional
relation; ignoring that relation would give an incorrect kernel.

Set $z=x_{c-1}$ and $x_0=0$. The coefficient of $f_k$ in
$b^{-1}L\Gamma$ is
\[
       ux_k+x_{k-1}-c_{c-k}(N)z\qquad(1\le k\le c-1).
\]
The equations for $k=c-1,\ldots,2$ determine every coefficient:
\[
 x_{c-1-j}=\sum_{k=0}^{j}c_k(N)(-u)^{j-k}z
                    \qquad(0\le j\le c-2).
\]
The equation for $k=1$ is exactly $q_N(-u)z=0$.
Conversely these coefficients give a primitive class whenever
that last equation holds. Uniqueness of the blow-up decomposition
makes extraction of $z$ inverse to this construction. C.1j now
identifies the kernel with the positivity radical. This proves
\eqref{eq:C1q-kernel}, including the cases of negative or zero
controlling degree, where the relevant Griffiths group vanishes.
\end{proof}

There is a genuine ample test, not just a formal exceptional
class. Suppose $Z$ is the smooth transverse base locus of two
sections of an ample line bundle $H$ on $M$, so $c=2$ and
$N=H|_Z\oplus H|_Z$. The pencil defines $f:X\to\mathbb P^1$,
and its closed graph realization gives
\[
 L=2\pi^*H-E=\pi^*H+f^*\mathcal O_{\mathbb P^1}(1)
                         \quad\text{ample}.
\]
Indeed this is the restriction of the ample external product
on $M\times\mathbb P^1$ to the closed graph. Here $u=2H|_Z$
and $q_N(-u)=c_1(N)-u=0$. Consequently
\begin{equation}\label{eq:C1q-pencil}
       \operatorname{rad}(q_{X,L})\simeq G^{m-1}(Z).
\end{equation}
A nonzero middle Griffiths class on such a pencil base would
therefore refute universal $B$. This is a criterion, not an
assertion that such a class has been constructed. If $M$ has
Tate cohomology, weak Lefschetz gives
$H^{2m-3}(Z,\mathbb Q)=0$; the required class would have zero
ordinary Abel--Jacobi invariant. The above-middle Nori and
quintic classes already established in this paper do not
supply this middle class.

For completeness, we can rule out a substantial stabilized
quintic attempt with its actual ample classes. Let $V\subset
\mathbb P^4$ be any smooth quintic threefold, let $r\ge0$,
$c\ge2$, and assume $r+c+3=2m$. Set
\[
 M=\mathbb P^4\times\mathbb P^{r+c-1},\qquad
 Z=V\times\mathbb P^r,
\]
using a linear inclusion for the second factor. Let $H,t$
be the two hyperplane classes. Every ample rational class on
$X=\operatorname{Bl}_Z(M)$ has the form
$L=\alpha\pi^*H+\beta\pi^*t-bE$ with $b>0$.
It also has $\beta>b$: take a line in a fibre over a point
of $V$ meeting the linear $\mathbb P^r$ once, without being
contained in it. Its strict transform has $L$-degree $\beta-b$.

The only possibly nonzero Griffiths group of $V$ is $G^2(V)$,
and $H G^2(V)=0$ because $G^3(V)=0$. Thus the controlling
group in \eqref{eq:C1q-kernel} is zero unless
\[
 e=\frac{r-c+1}{2}\ge0,
 \quad\text{in which case it is }G^2(V)t^e.
\]
Here
$N=\mathcal O_V(5)\boxplus
\mathcal O_{\mathbb P^r}(1)^{\oplus(c-1)}$.
On the Griffiths summand the terms containing $H$ vanish,
and the binomial theorem gives
\[
 q_N(-u)(g t^e)
      =(1-\beta/b)^{c-1}g t^{e+c-1}.
\]
Since $e\ge0$ implies $e+c-1\le r$, the last projective
coefficient is a genuine direct summand. Its coefficient
is nonzero because $\beta>b$. Therefore this map is
injective, and every such blow-up satisfies its own strict
middle positivity assertion. The apparent kernel at
$\beta=b$ lies outside the ample cone. This explicitly
records why this blow-up attempt does not refute universal
$B$; it neither proves that universal assertion nor rules
out other centers satisfying \eqref{eq:C1q-pencil}.
